% ── Chapter 6: Testing, Validation & Results ─────────────────────────────────
\chapter{Testing, Validation and Results}

\section*{Introduction}
\addcontentsline{toc}{section}{Introduction}

Chapters~4 and~5 designed LOGIQ and then built it. A decision-support system, however, is only as
valuable as the confidence its users place in it: it must be \emph{shown} to work. This chapter
validates the platform end to end and presents the results that justify the project.

Because the production source systems could not be connected directly during development, the validation
rests on a \textbf{faithful test environment} that exercises the entire chain --- from the sources,
through the ETL pipeline, the warehouse, and the backend API, up to the dashboards and alerts. The
guiding property of this environment is that the pipeline connects to it \emph{exactly as it would to
the real systems}: validating against the test environment therefore genuinely validates the production
design.

The chapter is organised in three movements. It first presents the \textbf{test apparatus} --- the
simulated source systems (\S\ref{sec:ch6-sources}), the business-realistic data generation
(\S\ref{sec:ch6-data}), its production scale (\S\ref{sec:ch6-scale}), and the automated data-quality
gate that admits only clean datasets (\S\ref{sec:ch6-quality}). It then reports the \textbf{end-to-end
validation} of the full chain (\S\ref{sec:ch6-e2e}). Finally, it presents the \textbf{results} the
platform delivers against the four objectives and the two business axes (\S\ref{sec:ch6-results}), and
closes on an honest \textbf{discussion} of what was achieved and where the work stops
(\S\ref{sec:ch6-discussion}).

% ═════════════════════════════════════════════════════════════════════════════
\section{Simulated Source Systems}
\label{sec:ch6-sources}

At the head of the test environment is a dedicated service --- built with \textbf{FastAPI} over
\textbf{PostgreSQL} --- that \textbf{simulates the company's five internal source systems}. Each system
is exposed behind a REST API, and the data it serves stands in for the operational applications scoped in
\S\ref{sec:ch4-dw}. Together the five cover every feed the warehouse needs (Table~\ref{tab:ch6-sources}).

% [VISUAL-IMPLEMENT] the five simulated source systems and the domain each provides
\begin{table}[H]
\centering
\renewcommand{\arraystretch}{1.35}
\rowcolors{2}{white}{tblrowalt}
\begin{tabularx}{\linewidth}{|>{\raggedright\arraybackslash}p{3.6cm}|>{\raggedright\arraybackslash}X|}
\hline
\rowcolor{tblheader}
\textcolor{white}{\textbf{Simulated system}} & \textcolor{white}{\textbf{Domain it provides}} \\
\hline
Core logistics system & Geography, delivery centres, the pricing grid, and the full parcel-event
history. \\
HR system & Companies, agencies, employees, and the organisational hierarchy. \\
Cash-box system & Operating expenses, freelance-driver payments, reimbursements, and internal
transfers. \\
Payroll system & Monthly payslips. \\
Transport system & Dedicated B2B transport requests and their stops. \\
\hline
\end{tabularx}
\caption{The five simulated source systems and the business domain each one provides to the warehouse.}
\label{tab:ch6-sources}
\end{table}

What makes the simulation valuable is not the data alone but the \textbf{fidelity of the contract}. Each
mock API reproduces the same \emph{endpoints, response shapes, authentication, pagination, and quirks}
as its real counterpart. Because the ETL pipeline depends only on this published contract --- never on a
system's internal database (the source decoupling of \S\ref{sec:ch4-etl}, TS17) --- it connects to the
simulation \textbf{exactly as it would to the production systems}, with no change to its extraction code.
Validating the platform against the simulation therefore genuinely validates the production design: the
same pipeline that passes here would run unaltered against the live systems.

% ═════════════════════════════════════════════════════════════════════════════
\section{Realistic Data Generation}
\label{sec:ch6-data}

The simulated systems are only as useful as the data they serve. That data is \textbf{synthetically
generated but business-realistic}, and it is precisely this realism --- rather than toy data --- that
makes the validation meaningful.

The generator works in \textbf{two tiers}. First, \emph{real reference data} seeds the dimensions: the
actual Algerian administrative geography (wilayas and communes), the real network of agencies and
delivery centres, the real company structure, and the real pricing grid. Second, \emph{synthetic
transactional data} is generated on top --- parcels and their lifecycles, payroll, operating expenses,
and transport requests --- so that the facts are plausible but the descriptive context is genuine.

Crucially, the synthetic data is shaped by \textbf{domain rules} drawn from how the business actually
behaves:

\begin{itemize}[noitemsep]
    \item \textbf{Working calendar} --- the Algerian week, with Friday as the only day off and reduced
          activity at weekends; no operational events fall on the non-working day.
    \item \textbf{Seasonality} --- the yearly rhythm of Algerian commerce: the Ramadan slowdown and the
          Eid rebound, the Black Friday surge, and the end-of-year peak (Figure~\ref{fig:ch6-seasonal}).
    \item \textbf{Geographic realism} --- parcel volume distributed across the wilayas in proportions
          that mirror the real activity, concentrated on the densest regions.
    \item \textbf{Labour-law payroll} --- salaries computed following Algerian labour-law rates, so
          payroll costs are realistic.
    \item \textbf{Full parcel lifecycle} --- each parcel passes through its complete multi-status
          lifecycle, not a single end state.
    \item \textbf{Internal consistency} --- costs, billing, and operations are kept mutually consistent,
          so financial and operational figures reconcile.
\end{itemize}

% [VISUAL-IMPLEMENT] illustrative seasonal parcel-volume chart (NOT measured data)
\begin{figure}[H]
\centering
\begin{tikzpicture}
\begin{axis}[
    width=\textwidth, height=6cm,
    ybar, bar width=8pt,
    ymin=0,
    symbolic x coords={Jan,Feb,Mar,Apr,May,Jun,Jul,Aug,Sep,Oct,Nov,Dec},
    xtick=data,
    x tick label style={font=\scriptsize},
    ylabel={Relative volume (illustrative)},
    ylabel style={font=\scriptsize},
    ytick style={draw=none},
    ymajorgrids, grid style={figslate!30, dotted},
    enlarge x limits=0.05,
    axis lines=left,
]
\addplot[fill=figbluetint, draw=figblue] coordinates
  {(Jan,90)(Feb,95)(Mar,105)(Apr,80)(May,120)(Jun,110)(Jul,100)(Aug,95)(Sep,105)(Oct,115)(Nov,160)(Dec,150)};
\node[font=\tiny, text=figslate, anchor=south] at (axis cs:Apr,80) {Ramadan};
\node[font=\tiny\bfseries, text=figorange, anchor=south] at (axis cs:Nov,160) {Black Friday};
\end{axis}
\end{tikzpicture}
\caption{Illustrative seasonal pattern of parcel volume produced by the generator --- the Ramadan
slowdown, the Black Friday surge, and the year-end peak. The values are \emph{illustrative}, shown to
convey the shape of the seasonality, not measured figures.}
\label{fig:ch6-seasonal}
\end{figure}

Taken together, these rules mean the test dataset does not merely fill the warehouse with rows: it
reproduces the \emph{behaviour} of the real activity, so that the dashboards, KPIs, and alerts are
exercised against patterns they will genuinely have to handle in production.

% ═════════════════════════════════════════════════════════════════════════════
\section{Scale of the Test Dataset}
\label{sec:ch6-scale}

A realistic \emph{shape} is not enough; the data must also have a realistic \emph{size}. The generator
produces the dataset at \textbf{production scale}, so the platform can be tested for performance and
correctness under genuine load rather than on a small sample (Table~\ref{tab:ch6-scale}).

% [VISUAL-IMPLEMENT] qualitative scale of the test dataset (exact counts deferred)
\begin{table}[H]
\centering
\renewcommand{\arraystretch}{1.35}
\rowcolors{2}{white}{tblrowalt}
\begin{tabularx}{\linewidth}{|>{\raggedright\arraybackslash}p{4.6cm}|>{\raggedright\arraybackslash}X|}
\hline
\rowcolor{tblheader}
\textcolor{white}{\textbf{Data domain}} & \textcolor{white}{\textbf{Approximate scale}} \\
\hline
Parcels generated & On the order of tens of thousands per day, over a multi-year horizon. \\
% TODO(achraf): confirm exact count (~220M staging events, ~27M parcels) from docs/dw-doc.md
Parcel-status events & The largest fact --- on the order of \emph{hundreds of millions} of rows. \\
% TODO(achraf): confirm exact payroll count from docs/mock-data-doc.md
Payroll & Tens of thousands of employee-months of payslips. \\
% TODO(achraf): confirm exact expense/transport counts
Expenses \& transport & Matching operating-expense and transport-request records over the same
horizon. \\
\hline
\end{tabularx}
\caption{The approximate scale of the generated test dataset. Figures are qualitative; exact counts are
to be confirmed from the design documents.}
\label{tab:ch6-scale}
\end{table}

This scale is what makes the validation credible. It is precisely the volume --- hundreds of millions of
parcel-status events --- that justifies the two design choices made earlier for performance: the
pre-computed \textbf{aggregate layer} (\S\ref{sec:ch4-dw}) and the \textbf{incremental} ETL load
(\S\ref{sec:ch5-data}). On a toy dataset neither could be meaningfully exercised; at production scale,
both are put under the same pressure they will face in operation.

% ═════════════════════════════════════════════════════════════════════════════
\section{Data-Quality Validation}
\label{sec:ch6-quality}

Realistic, production-scale data is worthless for validation if it is internally inconsistent: an
anomaly seen later in a dashboard could then be blamed on the data rather than the platform. To rule
this out, an automated \textbf{validation suite} of referential-integrity and business-rule checks runs
\emph{after every generation}. Only when \emph{all} checks pass is the dataset admitted to the ETL ---
establishing a clean, trustworthy baseline against which the rest of the platform can be tested
(Table~\ref{tab:ch6-quality}).

% [VISUAL-IMPLEMENT] data-quality check categories and what each verifies
\begin{table}[H]
\centering
\renewcommand{\arraystretch}{1.35}
\rowcolors{2}{white}{tblrowalt}
\begin{tabularx}{\linewidth}{|>{\raggedright\arraybackslash}p{3.9cm}|>{\raggedright\arraybackslash}X|}
\hline
\rowcolor{tblheader}
\textcolor{white}{\textbf{Check category}} & \textcolor{white}{\textbf{What it verifies}} \\
\hline
Referential integrity & Every cross-system reference resolves --- each foreign reference points to an
entity that actually exists. \\
Financial identities & Costs, billing, and reimbursement figures reconcile, so the financial picture is
self-consistent. \\
Distance identities & Transport distances are coherent --- for example, a request's leg distances sum to
its total. \\
Working calendar & No operational event falls on the non-working day (Friday). \\
Identifier formats & Identifiers follow their expected formats and conventions. \\
Test-entity isolation & Excluded test entities never leak into the generated data. \\
\hline
\end{tabularx}
\caption{The categories of data-quality check run after every generation; the dataset is admitted to the
ETL only when all of them pass.}
\label{tab:ch6-quality}
\end{table}

This quality gate is what makes the end-to-end validation that follows trustworthy. Because the input
data is guaranteed consistent before the pipeline ever touches it, any discrepancy observed downstream
--- in the warehouse, the API, or the dashboards --- can be attributed to the platform under test, not
to a flaw in its input.

% ═════════════════════════════════════════════════════════════════════════════
\section{End-to-End Validation}
\label{sec:ch6-e2e}

With a clean, production-scale dataset in place, the platform was exercised \textbf{end to end}: the
whole chain --- from the simulated sources, through the ETL pipeline, the warehouse, and the backend
API, to the dashboards and alerts --- was run as one, and the output of each stage was checked against
expectations (Figure~\ref{fig:ch6-e2e}).

% [VISUAL-IMPLEMENT] end-to-end validation chain (architecture style)
\begin{figure}[H]
\centering
\resizebox{\textwidth}{!}{%
\begin{tikzpicture}[
  >={Stealth[length=2.6mm]},
  stage/.style={draw, rounded corners=3pt, align=center, minimum width=2.2cm,
                minimum height=1.3cm, font=\small\bfseries, text=black!80},
  flow/.style={->, draw=figslate, line width=1pt}]

  \node[stage, draw=figslate, fill=black!4] (src) {Simulated\\sources};
  \node[stage, draw=figgreen,  fill=figgreentint, right=0.75cm of src] (etl) {ETL\\pipeline};
  \node[stage, draw=figteal,   fill=figtealtint,  right=0.75cm of etl] (wh)  {Warehouse};
  \node[stage, draw=figblue,   fill=figbluetint,  right=0.75cm of wh]  (api) {Backend\\API};
  \node[stage, draw=figblue,   fill=figbluetint,  right=0.75cm of api] (dsh) {Dashboards};
  \node[stage, draw=figorange, fill=figambertint, right=0.75cm of dsh] (alr) {Alerts};

  \draw[flow] (src) -- (etl);
  \draw[flow] (etl) -- (wh);
  \draw[flow] (wh)  -- (api);
  \draw[flow] (api) -- (dsh);
  \draw[flow] (dsh) -- (alr);
\end{tikzpicture}%
}
\caption{The end-to-end validation chain: the complete flow from the simulated sources to the alerts was
exercised as one, and each stage's output verified against expectations.}
\label{fig:ch6-e2e}
\end{figure}

Running the chain confirmed that the parts designed and built separately work \emph{together}: the ETL
extracts from the mock APIs and populates the warehouse layers; the aggregate layer feeds the backend's
analytics endpoints; the dashboards render the resulting KPIs; and the alerting service raises and
delivers notifications when a generated metric crosses a rule's threshold. Because the pipeline connects
to the simulation exactly as it would to production (\S\ref{sec:ch6-sources}), passing this chain
validates the production design itself, not merely a test fixture.

One behaviour was verified \textbf{independently} of the rest: the frontend's
\textbf{graceful degradation} (\S\ref{sec:ch5-backend}). By exercising the interface while the warehouse
was deliberately made unavailable, it was confirmed that the analytics endpoints fail safe and the
dashboard falls back to representative figures rather than breaking --- a property that holds, and is
demonstrable, even without a live warehouse.

% ═════════════════════════════════════════════════════════════════════════════
\section{Results}
\label{sec:ch6-results}

Exercised end to end against the test environment, the platform delivers concrete results against the
\textbf{four objectives} set in Chapter~3 and across the \textbf{two business axes}
(Table~\ref{tab:ch6-results}).

% [VISUAL-IMPLEMENT] the four objectives mapped to the results delivered
\begin{table}[H]
\centering
\renewcommand{\arraystretch}{1.35}
\rowcolors{2}{white}{tblrowalt}
\begin{tabularx}{\linewidth}{|>{\raggedright\arraybackslash}p{4.0cm}|>{\raggedright\arraybackslash}X|}
\hline
\rowcolor{tblheader}
\textcolor{white}{\textbf{Objective}} & \textcolor{white}{\textbf{Result delivered}} \\
\hline
Centralised warehouse & The data scattered across the source systems is consolidated into one
dimensional base; the figures are consistent across every view. \\
Reliable, automated ETL & The warehouse is refreshed automatically --- incrementally, idempotently, and
observably --- with no manual intervention. \\
Interactive dashboards & Both axes are navigable through their three views (Operations, Cost \&
Profitability, Performance), plus a consolidated overview, with KPI cards, charts, maps, and shared
filters. \\
Proactive alerting & Threshold rules are evaluated on schedule and raise multi-channel alerts the moment
a metric drifts beyond its acceptable level. \\
\hline
\end{tabularx}
\caption{The four objectives of the project and the results the validated platform delivers against
each.}
\label{tab:ch6-results}
\end{table}

Concretely, each axis page answers the decision questions it was designed for: \emph{Operations} (how
much volume, completion and delivery rates, geographic concentration), \emph{Cost \& Profitability}
(service cost, gross margin, payroll and expense weight), and \emph{Performance} (delays, attempts,
satisfaction signals). Populated with the generated data, the dashboards render these as live KPIs
(Figure~\ref{fig:ch6-result-kpi}), and the alerting service is observable in action when a generated
metric breaches a rule (Figure~\ref{fig:ch6-result-alert}).

% [VISUAL-PLACEHOLDER] results screenshot — a populated KPI dashboard page
\begin{figure}[H]
\centering
\visualplaceholder{A populated dashboard view --- KPI cards, charts, and a map filled with the generated
data for one axis view.}{Screenshot the running platform on a populated KPI page.}{A populated KPI
dashboard page, rendered from the generated test data.}
\label{fig:ch6-result-kpi}
\end{figure}

% [VISUAL-PLACEHOLDER] results screenshot — an alert firing
\begin{figure}[H]
\centering
\visualplaceholder{A threshold alert firing --- the raised alert and its in-app notification when a
metric crosses its rule.}{Screenshot the running platform when an alert is raised.}{A proactive alert
firing when a monitored metric crosses its threshold.}
\label{fig:ch6-result-alert}
\end{figure}

On the \textbf{performance} side, the design choices proved out under realistic load: because the
analytics endpoints read the pre-computed aggregate layer rather than scanning the fact tables, the
dashboards return within seconds even over the production-scale dataset, and the incremental refresh
keeps the warehouse current without a full reload.
% TODO(achraf): add measured numbers — dashboard response times, ETL refresh duration — once benchmarked
The two business axes are both delivered \textbf{in full}, each on its own financial perimeter, exactly
as scoped.

% ═════════════════════════════════════════════════════════════════════════════
\section{Discussion}
\label{sec:ch6-discussion}
% TODO

% ═════════════════════════════════════════════════════════════════════════════
\section*{Conclusion}
\addcontentsline{toc}{section}{Conclusion}
% TODO
